\documentclass[12pt]{article}
\usepackage{amsmath, amssymb}
\usepackage[margin=1in]{geometry}
\setlength{\parskip}{6pt}
\title{QUBO Formulation: Dynamic Spectrum Allocation}
\date{}
\begin{document}
\maketitle

\subsection*{Dynamic Spectrum Allocation}

\paragraph{Problem Statement.}
Given a set of radio links and a set of available communication channels, the objective is to assign exactly one channel to each link. Certain pairs of links experience mutual interference if they are assigned the same channel. Each link--channel assignment carries an associated utility score. The goal is to determine a channel assignment that maximizes the total utility across all links while penalizing interference between conflicting links.

\paragraph{Decision Variables.}
Let $L$ denote the set of radio links with cardinality $N_L = |L|$, and let $C$ denote the set of available channels with cardinality $N_C = |C|$. For each link $l \in L$ and channel $c \in C$, we introduce a binary decision variable $x_{l,c} \in \{0,1\}$. The variable $x_{l,c}$ equals $1$ if and only if link $l$ is assigned channel $c$, and $0$ otherwise.

\paragraph{QUBO Derivation.}
We convert the maximization problem into a minimization problem suitable for a Quadratic Unconstrained Binary Optimization (QUBO) formulation. Let $I \subseteq L \times L$ be the set of interfering link pairs, and let $P_{l_1, l_2}$ denote the interference penalty weight for pair $(l_1, l_2) \in I$. Let $U_{l,c}$ denote the utility score for assigning channel $c$ to link $l$.

\textbf{Step 1.} The original objective is to maximize $\sum_{l \in L} \sum_{c \in C} U_{l,c} x_{l,c} - \sum_{(l_1, l_2) \in I} P_{l_1, l_2} \sum_{c \in C} x_{l_1,c} x_{l_2,c}$. Flipping the sign to obtain a minimization objective yields:
\begin{align}
H_{\text{obj}}(x) = \sum_{l \in L} \sum_{c \in C} (-U_{l,c}) x_{l,c} + \sum_{(l_1, l_2) \in I} P_{l_1, l_2} \sum_{c \in C} x_{l_1,c} x_{l_2,c}.
\end{align}

\textbf{Step 2.} The hard constraint requires that each link receives exactly one channel. For every link $l \in L$, this is expressed as:
\begin{align}
\sum_{c \in C} x_{l,c} = 1.
\end{align}

\textbf{Step 3.} We enforce the constraints using quadratic penalty terms with weight $\lambda > 0$. The penalty Hamiltonian is:
\begin{align}
H_{\text{pen}}(x) = \sum_{l \in L} \lambda \left( \sum_{c \in C} x_{l,c} - 1 \right)^2.
\end{align}
Expanding the square for a fixed link $l$:
\begin{align}
\left( \sum_{c \in C} x_{l,c} - 1 \right)^2 &= \sum_{c \in C} x_{l,c}^2 + \sum_{c \neq c' \in C} x_{l,c} x_{l,c'} - 2 \sum_{c \in C} x_{l,c} + 1.
\end{align}
Applying the idempotent property $x_{l,c}^2 = x_{l,c}$ for binary variables:
\begin{align}
\left( \sum_{c \in C} x_{l,c} - 1 \right)^2 &= \sum_{c \in C} x_{l,c} + \sum_{c \neq c' \in C} x_{l,c} x_{l,c'} - 2 \sum_{c \in C} x_{l,c} + 1 \\
&= -\sum_{c \in C} x_{l,c} + \sum_{c \neq c' \in C} x_{l,c} x_{l,c'} + 1.
\end{align}
Multiplying by $\lambda$ and summing over all links gives the expanded penalty contribution:
\begin{align}
H_{\text{pen}}(x) = \sum_{l \in L} \left[ -\lambda \sum_{c \in C} x_{l,c} + \lambda \sum_{c \neq c' \in C} x_{l,c} x_{l,c'} + \lambda \right].
\end{align}

\textbf{Step 4.} The total QUBO Hamiltonian $H(x)$ is the sum of the objective and penalty terms:
\begin{align}
H(x) = H_{\text{obj}}(x) + H_{\text{pen}}(x).
\end{align}
Grouping linear terms, quadratic terms, and constants:
\begin{align}
H(x) = \sum_{l \in L} \sum_{c \in C} (-U_{l,c} - \lambda) x_{l,c} + \sum_{(l_1, l_2) \in I} P_{l_1, l_2} \sum_{c \in C} x_{l_1,c} x_{l_2,c} + \sum_{l \in L} \lambda \sum_{c \neq c' \in C} x_{l,c} x_{l,c'} + N_L \lambda.
\end{align}

\textbf{Step 5.} We map each pair $(l,c)$ to a single index $k = l N_C + c \in \{0, 1, \dots, N_L N_C - 1\}$. Reading off the coefficients from the expanded Hamiltonian, the QUBO matrix entries $Q_{k,k'}$ and offset $c$ are given in closed form by:
\begin{align}
Q_{k,k} &= -U_{l,c} - \lambda, \\
Q_{k,k'} &= \begin{cases}
P_{l_1, l_2} & \text{if } l_1 \neq l_2, (l_1, l_2) \in I, \text{ and } c = c', \\
\lambda & \text{if } l_1 = l_2 \text{ and } c \neq c', \\
0 & \text{otherwise},
\end{cases} \\
c &= N_L \lambda,
\end{align}
where $k = l N_C + c$ and $k' = l' N_C + c'$. The QUBO is then expressed as $H(x) = \sum_{k} Q_{k,k} x_k + \sum_{k < k'} Q_{k,k'} x_k x_{k'} + c$.

\paragraph{Penalty Weight Justification.}
The penalty weight $\lambda$ must be chosen sufficiently large to ensure that any constraint violation is strictly more costly than the maximum possible gain in the objective function. The maximum utility gain achievable by assigning an invalid number of channels to a link is bounded by $\max_{l \in L, c \in C} U_{l,c}$. Therefore, setting $\lambda > \max_{l \in L, c \in C} U_{l,c}$ guarantees that the energy penalty for violating the one-channel-per-link constraint dominates any potential utility improvement, forcing the optimizer toward feasible assignments.

\paragraph{Correctness Argument.}
Minimizing $H(x)$ over $\{0,1\}^{N_L N_C}$ recovers the optimal channel assignment because the Hamiltonian is constructed such that (i) all feasible solutions (exactly one channel per link) incur zero penalty and evaluate to the negative of the original objective function, and (ii) any infeasible solution incurs a penalty strictly greater than the maximum possible objective value. Consequently, the global minimum of $H(x)$ must correspond to a feasible assignment that minimizes the penalized energy, which is equivalent to maximizing the original utility minus interference penalties.

\paragraph{Illustrative Example.}
Consider a small instance with $N_L = 2$ links and $N_C = 2$ channels. Let $L = \{0, 1\}$ and $C = \{0, 1\}$. Suppose links $0$ and $1$ interfere, so $I = \{(0,1)\}$ with penalty $P_{0,1} = P$. Let the utility scores be $U_{0,0} = 10$, $U_{0,1} = 3$, $U_{1,0} = 4$, and $U_{1,1} = 9$. The binary variables are $x_{0,0}, x_{0,1}, x_{1,0}, x_{1,1}$.

Substituting these parameters into the general formulas yields the concrete Hamiltonian:
\begin{align}
H(x) &= (-10-\lambda)x_{0,0} + (-3-\lambda)x_{0,1} + (-4-\lambda)x_{1,0} + (-9-\lambda)x_{1,1} \\
&\quad + P x_{0,0} x_{1,0} + P x_{0,1} x_{1,1} + \lambda x_{0,0} x_{0,1} + \lambda x_{1,0} x_{1,1} + 2\lambda.
\end{align}
The constraints force $x_{0,0} + x_{0,1} = 1$ and $x_{1,0} + x_{1,1} = 1$. Evaluating the four feasible assignments:
\begin{itemize}
    \item $(x_{0,0}, x_{0,1}, x_{1,0}, x_{1,1}) = (1,0,1,0)$: Utility $10+4=14$, interference $P$. Energy $-14 + P$.
    \item $(1,0,0,1)$: Utility $10+9=19$, interference $0$. Energy $-19$.
    \item $(0,1,1,0)$: Utility $3+4=7$, interference $0$. Energy $-7$.
    \item $(0,1,0,1)$: Utility $3+9=12$, interference $P$. Energy $-12 + P$.
\end{itemize}
Assuming $\lambda > 10$ and $P < 5$, the assignment $(1,0,0,1)$ yields the lowest energy $-19$, corresponding to link $0$ on channel $0$ and link $1$ on channel $1$, which maximizes total utility while avoiding interference.
\end{document}